% -*- root: ../thesis.tex -*-
%!TEX root = ../thesis.tex
% this file is called up by thesis.tex
% content in this file will be fed into the main document

%level followed %by section, subsection


% ----------------------- paths to graphics ------------------------

% change according to folder and file names
\graphicspath{{chapter_PPT_semantization/figures/}}
% ----------------------- contents from here ------------------------

\section{Context within thesis}
This chapter studies how episodic memories change under systems consolidation in the parallel pathway theory (PPT, introduced in chapter \ref{ch:ppt}). It shows how memory transfer in the PPT can lead to the semantization of memories, and suggests that a hierarchy of parallel pathways from hippocampus to cortex could lead to a spatial gradient of memory semantization. The content of this chapter is based on my master's thesis \cite{Alevi2019,Alevi2019a}, which was carried out under supervision of Henning Sprekeler.

%\section{Contribution}
%
%\section{Introduction}
%
%\section{Results}
%
%\section{Methods}
%
%\section{Discussion}

\section{Chapter summary}

This chapter studied the semantization of episodic memories in the parallel pathway theory (PPT, chapter \ref{ch:ppt}) and hypothesised that a hierarchy of parallel pathways from hippocampus to cortex could lead to a spatial gradient of memory semantization and generalization. In the next chapter, we will explore how such different levels of semantic generalizations can be read out in a Bayes-optimal fashion, and how that can implement Bayesian change point detection in dynamic environments.

% ---------------------------------------------------------------------------
%: ----------------------- end of thesis sub-document ------------------------
% ---------------------------------------------------------------------------

